\originalchapter{电流与直流电路}\label{ch:dc}
\originalsection{漂移、密度与电阻}
电荷持续穿过截面形成电流$I=\dd Q/\dd t$, 单位安培(A). 约定方向是正电荷运动方向, 在金属中通常与电子漂移方向相反. 电流密度$\jj=nq\vct v_{\rm d}$为单位面积的有向电流, $I=\int\jj\cdot\dd\vct A$. 晶格碰撞使平均漂移速度很小, 电路开始响应的速度则由电磁场传播决定, 不是某个电子从电池跑到灯泡的速度.

电荷守恒在任意体积中为$\dd Q_{\rm in}/\dd t=-\oint\jj\cdot\dd\vct A$, 局部形式为$\partial_t\rho+\nabla\cdot\jj=0$. 稳态节点不积累电荷, 流入总电流等于流出总电流, 即基尔霍夫节点定律.

欧姆材料在给定温度等条件下近似满足$\jj=\sigma_{\rm c}\ee$, 电导率$\sigma_{\rm c}$不是面电荷密度. 均匀导线长$\ell$、面积$A$, 电阻率$\rho_{\rm r}=1/\sigma_{\rm c}$, 有$R=\rho_{\rm r}\ell/A$、$\Delta V=IR$. 欧姆定律是材料近似, 不是所有元件都必须遵守的定义. 电阻把场能转为热, 功率$P=I\Delta V=I^2R=\Delta V^2/R$.

\originalsection{电源与回路}
静电场沿闭合回路净作功为零, 电阻却不断耗能. 电池通过化学作用把电荷从较低势能处搬到较高势能处; 电动势$\mathcal E$定义为电源的非静电作用每单位电荷提供的能量, 虽称“力”却以伏特为单位. 有内阻$r$的电池放电时端电压为$\mathcal E-Ir$.

在回路尺度小于相关电磁波波长、且回路外磁通变化可忽略时, 各节点可赋电势, 一圈电势变化代数和为零, 即基尔霍夫回路定律. 顺电流经过电阻取$-IR$, 从电池负极到正极取$+\mathcal E$. 先选假定电流方向, 求出负值表示真实方向相反.

串联电阻电流相同, 势差相加, $R_{\rm eq}=\sum R_i$. 并联电阻势差相同, 电流相加, $1/R_{\rm eq}=\sum1/R_i$. 理想电压表内阻很大, 接在两点之间; 理想电流表内阻很小, 串入支路. 将理想电流表直接并在电源上会改变电路而造成很大电流, 并非无扰动测量.

节点是由理想导线连成、可赋同一电势的区域, 支路是两个节点之间含元件的一段. 当电路有多个电源和多个支路, 独立节点方程加独立回路方程组成线性系统; 同一回路重复写几次并不会增加独立信息.

\originalsection{电容暂态}
串联$R,C$接恒定$\mathcal E$, 以流向正极板的电流为正, $I=\dot q$, 回路式为$R\dot q+q/C=\mathcal E$. 分离变量或解一阶线性方程, 初始$q(0)=q_0$, 得
\[
q(t)=C\mathcal E+(q_0-C\mathcal E)e^{-t/(RC)},\quad I(t)=\frac{C\mathcal E-q_0}{RC}e^{-t/(RC)}.
\]
时间常数$\tau=RC$是偏离最终值缩至原来$1/e$所需的时间. 放电时$\mathcal E=0$, 得$q=q_0e^{-t/RC}$. 有限电流不能使$q$瞬间跳变, 所以电容电压连续. 只有初始未充电的电容在闭合的瞬间电压为零, 才可在求初始支路电流时当作短路; 稳态直流时它的电流为零, 可当作开路.

对复杂线性电阻网络, 从电容端看见的开路电压$V_{\rm th}$和撤掉独立理想电源后的等效电阻$R_{\rm th}$给出$R_{\rm th}\dot q+q/C=V_{\rm th}$. 将理想电压源置零意味着短接, 理想电流源置零意味着断开. 这使多支路RC题仍归为同一个一阶过程.


上述RC指数可以直接积分验证. 未充电时$\dot q=(C\mathcal E-q)/(RC)$, 分离变量:
\[
\int_0^{q(t)}\frac{\dd q'}{C\mathcal E-q'}=\int_0^t\frac{\dd t'}{RC},\quad
-\ln\frac{C\mathcal E-q(t)}{C\mathcal E}=\frac t{RC}.
\]
指数化便得$q=C\mathcal E(1-e^{-t/RC})$. 对任意$q_0$, 改为积分偏离量$q-C\mathcal E$, 它满足$\dot{(q-C\mathcal E)}=-(q-C\mathcal E)/(RC)$. 电阻始终消耗$I^2R$, 不能把电源给予的总能量全当作电容最终储能.
